\section*{\centering ABSTRACT}
\addcontentsline{toc}{section}{ABSTRACT}

Mobile service robots are typically delivered with a closed software ecosystem: a proprietary
application, no documentation, and no published communication protocol. This severely limits any
customization by the end user. The \textsc{Cybel} project, carried out as part of a final-year
internship at HESTIM, aims to design an independent control and interaction platform for a CIOT
TY1251D-03195 mobile reception robot, consisting of a ROS-based navigation chassis and an Android
\emph{« upper body »}.

The approach adopted relies on an incremental, non-destructive reverse engineering of the robot's
internal communication protocol: scanning of exposed network services, introspection of ROS
topics and services via \textbf{rosbridge}/\textbf{rosapi}, and systematic verification of the
actual effect of each command before considering it functional. On this basis, a three-layer
architecture was designed:
\begin{itemize}
    \item a Python SDK offering interchangeable simulated (\emph{mock}) and real implementations;
    \item a FastAPI backend (REST~+~WebSocket);
    \item a Vite/TypeScript operator web interface;
\end{itemize}
complemented by a visitor kiosk interface (\texttt{frontend-kiosk/}) and two lightweight native
Android applications (\texttt{CybelTTSBridge}, \texttt{CybelVisitorKiosk}).

At this stage of the internship, connectivity with the robot has been established; the
telemetry, speed control, and autonomous navigation protocols have been reverse-engineered and
integrated; text-to-speech (TTS) works through a dedicated Android application triggered via
\texttt{am broadcast}; a complete operator interface (SLAM map, LiDAR, detected visitors, tour
management, emergency stop during a visit) is operational; and an embedded deployment on Termux
(\emph{backend lite} + kiosk) has been validated on the tablet. The visitor interface has been
refocused on an autonomous guided tour of the laboratory: eight stops synchronized from
\texttt{data/knowledgeV2-lab.json} to \texttt{data/lab\_tour.json}, coordinate-based navigation
(\texttt{/navi\_goal}), and sequential voice announcements.

This report presents the context, problem statement, state of the art, design, methodology,
implementation carried out, results obtained, as well as a critical analysis of the project's
limitations and perspectives.

\bigskip
\noindent\textbf{Keywords:} service robotics, protocol reverse engineering, ROS, rosbridge,
FastAPI, human–robot interface, autonomous navigation, text-to-speech, Termux, Android WebView.